\section{DDR3 channels and controllers}
\label{sec:ddr}

\subsection{From the ports to the chips}
Each channel (\cref{fig:ddr}) is a LiteDRAM core: a \textbf{crossbar} that arbitrates between ports per bank, a
\textbf{controller} with one bank machine per DRAM bank, a refresher and a command multiplexer that enforces the
DDR3 timings (LiteDRAM's MT41K256M8 module, the same geometry and timings as the SK hynix H5TC2G83 chips on the card),
and the \textbf{PHY}. The controller runs at 1:4: one 125\,MHz sys cycle carries four DRAM clock cycles, so the native
port is 512 bits wide and the data rate is 1000\,MT/s on 64 bits, \SI{8}{GB/s} raw per channel. A BIST run after
calibration measures \SI{7.1}{GB/s} writing and \SI{7.2}{GB/s} reading on each channel.

Ports used in \code{GemmSoC}: the bridge has a write port and a read port on each channel; the kernel's A, B and C
masters each get a port on the channel of their matrix, in the accel domain (the crossbar inserts the asynchronous
FIFOs); the BIST generator and checker have their own ports. Channel 0 also keeps an idle Wishbone port from the
standard integration.

\begin{figure}[t]
  \centering
  \resizebox{\linewidth}{!}{% One DDR3 channel: crossbar ports, LiteDRAM controller, WL7DDRPHY, the fly-by DRAM and the host calibration loop.
\begin{tikzpicture}[x=1cm,y=1cm, every node/.append style={font=\sffamily\footnotesize}]
  \tikzset{w/.style={minimum width=#1}, port/.style={blk, minimum width=30mm, minimum height=6mm, font=\sffamily\scriptsize}}
  % ports
  \node[port, fill=dSys]   (p1) at (0, 1.8)  {DMA bridge write port (sys)};
  \node[port, fill=dSys]   (p2) at (0, 1.05) {DMA bridge read port (sys)};
  \node[port, fill=dAccel] (p3) at (0, 0.3)  {kernel port A/B/C (accel, CDC)};
  \node[port, fill=dSys]   (p4) at (0,-0.45) {BIST generator + checker};
  \node[port, fill=black!5, draw=black!40, text=black!60] (p5) at (0,-1.2) {CH0 only: Wishbone (idle, no CPU)};
  % crossbar
  \node[sys, w=20mm, minimum height=36mm] (xb) at (3.4,0.3) {crossbar\\[2pt]\scriptsize per-bank\\ \scriptsize arbitration\\[2pt]\scriptsize 512\,b data\\ \scriptsize 1:4 words};
  \foreach \p in {p1,p2,p3,p4,p5} { \draw[biln] (\p.east) -- (\p.east -| xb.west); }
  % controller
  \node[sys, w=28mm, minimum height=36mm, align=center] (mc) at (6.75,0.3) {};
  \node[font=\sffamily\footnotesize\bfseries, anchor=north] at (mc.north) {LiteDRAM controller};
  \foreach \i in {0,...,7} {
    \node[draw=black!50, fill=white, minimum width=2.4mm, minimum height=7mm, inner sep=0pt] (bm\i) at (5.73+\i*0.29, 0.62) {};
  }
  \node[font=\sffamily\scriptsize] at (6.75,1.4) {8 bank machines};
  \node[draw=black!50, fill=white, minimum width=24mm, font=\sffamily\scriptsize] (mux) at (6.75,-0.25) {command multiplexer};
  \node[draw=black!50, fill=white, minimum width=24mm, font=\sffamily\scriptsize] (rf)  at (6.75,-0.9) {refresher · timing checks};
  \node[note] at (6.75,-1.35) {MT41K256M8 timings, 1:4};
  \draw[wbi] (xb) -- (mc);
  % PHY
  \node[ddr, w=30mm, minimum height=36mm] (phy) at (10.55,0.3) {};
  \node[font=\sffamily\footnotesize\bfseries, anchor=north] at (phy.north) {\code{WL7DDRPHY}};
  \node[align=center, font=\sffamily\scriptsize] at (10.55,0.15)
       {OSERDES / ISERDES 8:1\\ (4 phases \x DDR)\\[3pt] read path: IDELAY per DQ,\\ read bitslip per lane\\[3pt]
        write: bitslip per lane\\[3pt] \textbf{CK phase}: channel MMCM,\\ steps of 1/56 VCO \(\approx\) 17.9\,ps};
  \draw[wbi] (mc) -- node[lbl, above] {DFI} (phy);
  % DRAM: DQ lanes point to point, CK/command/address fly-by
  \foreach \c in {0,...,7} {
    \node[draw=cPurple!70!black, fill=cPurple!12, minimum width=7mm, minimum height=9mm, font=\sffamily\scriptsize,
          inner sep=0pt] (d\c) at (13.0+\c*0.85, 0.3) {\(\times\)8};
    \draw[cPurple!70!black, thick] (d\c.north) -- (d\c.north |- 0,1.25);
    \node[font=\sffamily\tiny, text=cPurple!70!black, anchor=south] at (d\c.north |- 0,1.25) {DQ\c};
    \draw[cPurple!70!black] (d\c.south) -- (d\c.south |- 0,-0.55);
  }
  \node[draw=black!30, fill=black!4, dashed, minimum width=7mm, minimum height=9mm, font=\sffamily\scriptsize,
        text=black!45, inner sep=0pt] (d8) at (13.0+8*0.85,0.3) {ECC};
  \draw[line width=1.6pt, cPurple!70!black] (phy.east |- 0,1.25) -- (d7.north |- 0,1.25);
  \node[note, anchor=south] at (12.45,1.3) {8 DQ/DQS lanes};
  \draw[very thick, cPurple!70!black, -{Stealth}] (phy.east |- 0,-0.55) -- ([xshift=4mm]d7.south |- 0,-0.55)
        node[right, font=\sffamily\scriptsize, align=left, text=cPurple!70!black] {CK, cmd, addr\\(fly-by)};
  \node[note, anchor=north, text width=64mm] at (16.0,-0.85) {CK reaches the far chips later than the near ones.
        HR banks (no ODELAY, so no write leveling): one CK phase has to suit every lane.};
  % host calibration loop
  \node[hst, align=left, text width=118mm] (cal) at (6.6,-3.0)
       {\textbf{Host calibration} (\file{boards/dupcb/ddr\_init.py}, CSRs over PCIe), per channel after every load:\\[1pt]
        \scriptsize 1.~For each CK phase step over one tCK: DRAM init; per-lane write bitslip, read bitslip and read delay
        scan with random bursts; 64\,MB BIST. A step passes if every lane calibrates and the BIST is clean.\\
        \scriptsize 2.~CK to the centre of the longest passing run; calibrate again.\quad
        3.~200 random addresses + BIST (write once, read twice).};
  \draw[ctl] (phy.south |- cal.north) -- node[lblw, right] {MMCM step, DFI, delays} (phy.south);
  \draw[ctl] (cal.west) -| node[lblw, pos=0.75, left] {BIST} ([xshift=-4mm]p4.west) -- (p4.west);
\end{tikzpicture}
}
  \caption{One DDR3 channel. The PHY's CK phase is moved by the host through the channel's MMCM; the calibration
    loop runs on the host over PCIe after every bitstream load (and after the driver load).}
  \label{fig:ddr}
\end{figure}

\subsection{PHY and calibration}
The DDR banks of this FPGA are high-range (HR) banks: they have no output delay lines, so the PHY cannot move DQS per
lane to do write leveling. The chips sit on a fly-by bus, so the far chips see CK later than the near ones. With
LiteX's stock Artix-7-style PHY there was no write timing that worked for every lane under traffic. The project
uses \code{WL7DDRPHY} (\file{boards/dupcb/wl7ddrphy.py}), from a sister card's design: each channel has its own
MMCM, and the host moves CK against DQ/DQS in steps of 1/56 of the VCO period (\(\approx\)\SI{17.9}{ps}).

Calibration (\file{boards/dupcb/ddr\_init.py}) runs on the host through the CSRs. For every CK phase step across one
clock period it initialises the DRAM, searches each lane's write bitslip, read bitslip and read delay with random
bursts, and runs a 64\,MB BIST. A step passes only if every lane calibrates and the BIST is clean; single-burst checks
alone pass at phases where streaming traffic fails. CK then goes to the centre of the longest passing run, the lanes
are calibrated again, and a random-address check and a final BIST (write once, read twice) confirm the channel. The
measured write windows were 893\,ps on CH0 and 1429\,ps on CH1. A full calibration of both channels takes about 15\,s.

\subsection{Where the matrices are}
The build's \code{memory\_map} decides which channel each kernel master is wired to; the host must use the same map,
and the run scripts read it from the build directory. Addresses are per channel; there is no common address space.
\Cref{fig:memmap} shows the default map, which separates the read traffic of A from that of B. \code{gemm\_runner}
checks that every matrix, including the slack of whole 8\,KB buffers that the C download reads, ends before the next
region on its channel.

\begin{figure}[ht]
  \centering
  \resizebox{0.95\linewidth}{!}{% Where A, B and C live in the two DDR3 channels (default memory_map a=ch0, b=ch1, c=ch1).
\begin{tikzpicture}[x=1cm,y=1cm, every node/.append style={font=\sffamily\footnotesize}]
  \def\L{14}  % 2 GB = 14 cm
  \foreach \ch/\y in {CH0/0, CH1/-1.35} {
    \draw[black!60, fill=black!3] (0,\y-0.35) rectangle (\L,\y+0.35);
    \node[anchor=east, font=\sffamily\small\bfseries] at (-0.15,\y) {\ch};
  }
  % A on CH0 from 0; example 8192^3: 256 MB
  \fill[cBlue!35] (0,-0.35) rectangle (\L/8,0.35);
  \draw[cBlue!70!black, dashed] (\L/8,-0.35) -- (\L/8,0.35);
  \node at (\L/16,0) {A};
  \node[anchor=west, font=\sffamily\scriptsize, text=black!60] at (\L/8+0.1,0) {room for A up to the whole 2\,GB channel};
  % B on CH1 from 0, C on CH1 from 1 GB
  \fill[cOrange!45] (0,-1.7) rectangle (\L/8,-1.0);
  \node at (\L/16,-1.35) {B};
  \node[anchor=west, font=\sffamily\scriptsize, text=black!60] at (\L/8+0.1,-1.35) {B \(\le\) 1\,GB};
  \fill[cGreen!40] (\L/2,-1.7) rectangle (\L/2+\L/8,-1.0);
  \node at (\L/2+\L/16,-1.35) {C};
  \node[anchor=west, font=\sffamily\scriptsize, text=black!60] at (\L/2+\L/8+0.1,-1.35) {C \(\le\) 1\,GB, read back in whole 8\,KB buffers};
  \draw[black!70, thick] (\L/2,-1.75) -- (\L/2,-0.95);
  % address axis
  \foreach \a/\t in {0/0x0000\_0000, 0.25/0x2000\_0000, 0.5/0x4000\_0000, 0.75/0x6000\_0000, 1/0x8000\_0000} {
    \draw[black!50] (\a*\L,-1.8) -- (\a*\L,-1.95);
    \node[anchor=north, font=\ttfamily\scriptsize] at (\a*\L,-1.95) {\t};
  }
  \foreach \a/\t in {0.25/512\,MB, 0.5/1\,GB, 0.75/1.5\,GB, 1/2\,GB} {
    \node[anchor=south, font=\sffamily\scriptsize, text=black!55] at (\a*\L,0.38) {\t};
  }
  \node[note, anchor=north west, text width=13.5cm, align=left] at (0,-2.45)
    {Default \code{memory\_map}: A on CH0, B and C on CH1; shown for 8192\(^3\) (256\,MB each). When matrices share a
     channel, \code{gemm\_runner} puts B at 512\,MB and C at 1\,GB of that channel, and rejects matrices that would run
     into the next region. Byte addresses are per channel; there is no common address space across channels.};
\end{tikzpicture}
}
  \caption{Placement of A, B and C in the two DDR3 channels (default memory map).}
  \label{fig:memmap}
\end{figure}
